\documentclass[11pt]{article}
\usepackage[margin=1in]{geometry}
\usepackage{amsmath,amssymb,amsthm,mathtools}
\usepackage{graphicx}
\usepackage{xcolor}
\usepackage{microtype}
\usepackage[colorlinks=true,linkcolor=teal!70!black,citecolor=teal!70!black,urlcolor=teal!70!black]{hyperref}

\newtheorem{theorem}{Theorem}
\newtheorem{remark}[theorem]{Remark}
\newcommand{\R}{\mathbb R}
\newcommand{\Nrm}[1]{\left\lVert #1\right\rVert}

\title{A Sub-$p$ Fine Index Forbids the Sharp Uniform\\
Limiting Embedding}
\author{Run \texttt{fa\_banach\_001}, \texttt{agent\_lane\_14}}
\date{12 August 2026}

\begin{document}
\maketitle

\section*{Status and source question}

\textbf{Status:} candidate full negative answer to the sharp uniform
extension, likely valid pending expert review.

Kolyada~\cite{Kolyada} proves a sharp anisotropic Besov--Lorentz estimate when
every fine index $\theta_j\ge p$.  The factors
$(1-\beta_j)^{1/\theta_j}$ keep the constant bounded as the smoothness
parameters tend to one.  Immediately after Theorem~1.2, equation~(1.4), on
source PDF page~3, the paper says that the case $\theta_j<p$ for some $j$
remains open.

\begin{center}
\includegraphics[width=\textwidth]{figures/open_problem_crop.png}
\end{center}

The source also stresses that the same embedding \emph{without} these factors
is already known for all $1\le\theta_j\le\infty$.  Thus the nontrivial open
content is whether the sharp normalized estimate survives with a constant
bounded in the limiting regime.  The answer is no, even if only one fine index
is below $p$.

\section*{Proof intuition}

Let all smoothness indices equal $\rho$ and let $\rho\uparrow1$.  The normalized
directional Besov seminorms of a fixed $W^{1,p}$ function remain bounded; this
is the elementary upper half of the Bourgain--Brezis--Mironescu limit.  If one
fine index is $s<p$ and the others equal $p$, however, the target harmonic
fine index $\theta$ is strictly below $p$.  A uniform extension would therefore
force $W^{1,p}$ into the too-small Lorentz space $L^{p^*,\theta}$.

An explicit logarithmic radial singularity makes this obstruction quantitative.
It has an $L^p$ gradient because its logarithmic exponent is above $1/p$, but
its critical Lorentz norm diverges because the same exponent is below
$1/\theta$.  Before the limit, the gap
$1/q_\rho-1/p^*=(1-\rho)/n$ exponentially cuts off the logarithmic integral.
Its exact Laplace asymptotic gives the blow-up rate of the best constant.

\section*{Quantitative counterexample}

For $0<\rho<1$ and $1\le r<\infty$, write
\[
 [f]_{j,\rho,r}
 =\left(\int_0^\infty
       \big[h^{-\rho}\omega_j(f;h)_p\big]^r\,\frac{dh}{h}
  \right)^{1/r}.
\]

\begin{theorem}\label{thm:counterexample}
Let $n\ge2$, $1<p<n$, and $1\le s<p$.  Put
\[
 \theta_1=s,\qquad \theta_2=\cdots=\theta_n=p,
 \qquad
 \frac1\theta=\frac1n\left(\frac1s+\frac{n-1}{p}\right),
\]
and $q_\rho=np/(n-\rho p)$.  Let $C_\rho$ be the least constant for
which the proposed extension of source equation~(1.4), with
$\beta_1=\cdots=\beta_n=\rho$, could hold:
\begin{equation}\label{eq:proposed}
 \Nrm{f}_{L^{q_\rho,\theta}}
 \le C_\rho
 \prod_{j=1}^n
 \left[(1-\rho)^{1/\theta_j}[f]_{j,\rho,\theta_j}\right]^{1/n}.
\end{equation}
Set
\[
 \Gamma=\frac1\theta-\frac1p
 =\frac1n\left(\frac1s-\frac1p\right)>0.
\]
For every $0<\varepsilon<\Gamma$ there are $c_\varepsilon>0$ and
$\rho_\varepsilon<1$ such that
\[
 C_\rho\ge c_\varepsilon(1-\rho)^{-(\Gamma-\varepsilon)}
 \qquad(\rho_\varepsilon<\rho<1).
\]
In particular, no constant bounded as $\rho\uparrow1$ can extend the source's
sharp estimate when even one fine index is below $p$.
\end{theorem}

\begin{proof}
Fix $0<\varepsilon<\Gamma$ and set $a=1/p+\varepsilon$, so
$1/p<a<1/\theta$.  Let
\[
 \gamma=\frac{n-p}{p},\qquad L(r)=\log\frac er,
\]
and choose a radial cutoff $\eta\in C_c^\infty([0,\infty))$ which equals one
for all sufficiently small $r$ and vanishes for $r\ge e^{-2}$.  Define, with
an arbitrary value assigned at the origin,
\begin{equation}\label{eq:witness}
 f(x)=\eta(|x|)|x|^{-\gamma}L(|x|)^{-a}.
\end{equation}

Near the origin,
\[
 |f(x)|^p\asymp |x|^{-(n-p)}L(|x|)^{-ap},
 \qquad
 |\nabla f(x)|^p\asymp |x|^{-n}L(|x|)^{-ap}.
\]
The first function is locally integrable, and polar coordinates turn the
gradient integral into
\[
 \int_0^{e^{-3}}\frac{dr}{rL(r)^{ap}}<\infty
\]
because $ap>1$.  Hence $f\in W^{1,p}(\R^n)$.

For $0<h\le1$, the Sobolev translation inequality gives
\[
 \omega_j(f;h)_p\le h\Nrm{\partial_j f}_p,
\]
while for $h>1$ one has $\omega_j(f;h)_p\le2\Nrm f_p$.  Therefore, for
every finite $r\ge1$,
\begin{align}\label{eq:besovbound}
 (1-\rho)[f]_{j,\rho,r}^{,r}
 &\le (1-\rho)\Nrm{\partial_j f}_p^r
       \int_0^1h^{r(1-\rho)}\,\frac{dh}{h}
 \notag\\
 &\quad +(1-\rho)(2\Nrm f_p)^r
       \int_1^\infty h^{-r\rho}\,\frac{dh}{h}\\
 &=\frac{\Nrm{\partial_j f}_p^r}{r}
   +\frac{1-\rho}{r\rho}(2\Nrm f_p)^r.
\end{align}
Thus every factor on the right of \eqref{eq:proposed} stays bounded as
$\rho\uparrow1$.

It remains to estimate the left side.  Put $p^*=np/(n-p)$.  Since the
singular part of $f$ is radial decreasing, there are $c,C,t_0>0$ such that
\begin{equation}\label{eq:rearrangement}
 f^*(t)\ge c\,t^{-1/p^*}\left(\log\frac Ct\right)^{-a}
 \qquad(0<t<t_0).
\end{equation}
Moreover,
\[
 \frac1{q_\rho}-\frac1{p^*}=\frac{1-\rho}{n}.
\]
Using \eqref{eq:rearrangement} in the Lorentz norm and then setting
$u=\log(C/t)$ gives
\begin{align*}
 \Nrm f_{L^{q_\rho,\theta}}^\theta
 &\ge c\int_0^{t_0}
 t^{\theta(1-\rho)/n}
 \left(\log\frac Ct\right)^{-a\theta}\frac{dt}{t}\\
 &=c'\int_{u_0}^{\infty}
 e^{-\theta(1-\rho)u/n}u^{-a\theta}\,du.
\end{align*}
Because $a\theta<1$, the last integral is comparable, as $\rho\uparrow1$,
to $(1-\rho)^{a\theta-1}$.  Consequently
\[
 \Nrm f_{L^{q_\rho,\theta}}
 \ge c''(1-\rho)^{a-1/\theta}
 =c''(1-\rho)^{-(\Gamma-\varepsilon)}.
\]
Combining this with the uniform bound \eqref{eq:besovbound} proves the stated
lower bound for $C_\rho$.
\end{proof}

\begin{remark}[Concrete parameters]
One may take $(n,p,s)=(3,2,1)$.  Then $\theta=3/2$, $p^*=6$, and
$a=7/12$ works.  The witness is
$|x|^{-1/2}[\log(e/|x|)]^{-7/12}$ near the origin, its gradient is square
integrable, and the proposed constant grows at least like
$(1-\rho)^{-1/12}$.
\end{remark}

\section*{Verification, scope, and novelty}

The included script evaluates the incomplete-Gamma expression for the final
Laplace integral in the concrete example; it is not part of the proof.  The
argument only refutes uniform preservation of the sharp limiting factors.
For a fixed $\rho<1$, those factors are positive constants, and the older
unnormalized anisotropic embedding remains compatible with this result.

Run indexes and bounded arXiv searches through 12 August 2026 used the exact
title, arXiv id, author, quoted open phrase, and combinations of
``$\theta_j<p$'', anisotropic Besov, Lorentz endpoint, and sharp limiting
constant.  They found no later source recording this obstruction.  Novelty
confidence is moderate and requires specialist citation review.

\section*{Human-review recommendation}

First confirm the intended uniformity from the paragraphs surrounding source
equation~(1.4): without it, the question is already covered by the known
no-factor embedding.  Then check \eqref{eq:besovbound} and
\eqref{eq:rearrangement}.  These are the only substantive estimates; the
remaining blow-up is the elementary Gamma-integral scaling.

\begin{sloppypar}
\raggedright
\begin{thebibliography}{9}
\bibitem{Kolyada}
V.~I. Kolyada,
\emph{Rearrangement estimates and limiting embeddings for anisotropic Besov
spaces}, arXiv:2306.13938 (2023).
\end{thebibliography}
\end{sloppypar}

\end{document}
